\documentclass[11pt]{article}

\usepackage[T1]{fontenc}
\usepackage[margin=1in]{geometry}
\usepackage{graphicx}
\usepackage{amsmath}
\usepackage{booktabs}
\usepackage{array}
\usepackage{float}
\usepackage{xcolor}
\usepackage[hidelinks]{hyperref}

\emergencystretch=2em

\graphicspath{{figures/mc_bio_retraining/}{reports/figures/mc_bio_retraining/}}

\title{From Conditional Monte Carlo Diagnostics to FTV-Calibrated Consistent Rollout Models}
\author{}
\date{Progress update}

\begin{document}
\maketitle

\begin{abstract}
I used the conditional Monte Carlo experiment to stress-test the previously
successful consistent graph rollout model under the way I ultimately want to
use it: patient-level uncertainty estimation and treatment-response reasoning.
The original scheduled-sampling model produced coherent graph rollouts, but the
Monte Carlo analysis revealed a specific biological calibration problem. The
deterministic T0-to-T3 rollout underpredicted mean functional tumor volume
(FTV) by 12.51 mL, and the residual Monte Carlo correction overcompensated by
moving the Monte Carlo mean above the observed mean. Raw 90\% intervals
under-covered, and conformal intervals reached nominal coverage only after
substantial widening. These findings motivated a targeted retraining grid with
aggregate FTV and alive-burden objectives. The best retrained model reduced
T0-to-T3 FTV MAE from
12.88 mL to 5.88 mL while leaving spatial metrics essentially unchanged. I then
reran the same conditional Monte Carlo sampler on the retrained checkpoints.
For the primary retrained model, T0-to-T3 MC-mean FTV MAE dropped from
15.26 mL to 7.61 mL, raw 90\% coverage increased from 0.766 to 0.876, conformal
90\% interval width shrank from 59.51 mL to 27.95 mL, and MC alive-count error
dropped from 40.78 to 2.69 supervoxels. The interpretation is that the Monte
Carlo experiment did not invalidate the earlier rollout model; it identified
the next modeling requirement: endpoint-calibrated tumor burden before
uncertainty estimates can be trusted for response analysis. The conclusion from
this phase is that \texttt{bio\_ftv020\_alive005} is the best balanced
candidate to carry forward. Its predictions can now support patient-level
imaging questions that the original rollout model could not answer reliably:
what future MRI FTV is likely after conditioning on available visits, how wide
the plausible FTV range is, whether the observed endpoint falls inside a
calibrated interval, and how much active tumor burden remains in the graph.
That is meaningful because the model is no longer only a plausible spatial
forecaster; it is becoming an endpoint-calibrated uncertainty model for MRI
tumor-burden response analysis.
\end{abstract}

\section{Why this experiment mattered}

The previously retained model, \texttt{v2\_sched\_samp}, was already a strong
consistent graph rollout model. It used scheduled sampling and learned to roll
registered tumor graphs through the I-SPY2 imaging sequence:
\[
T0 \rightarrow T1 \rightarrow T2 \rightarrow T3.
\]
Each graph node represents a tumor supervoxel with a 3D position, imaging
features, a volume-like feature, and an alive probability. The model predicts
how those node states evolve across treatment.

That model was useful as a spatial rollout model, but I wanted to know whether
it could also serve as the center of a patient-level uncertainty distribution.
This is a harder question. A graph trajectory can look plausible while still
being wrong about total active tumor burden. For treatment-response modeling,
that distinction matters because FTV is the clinically interpretable endpoint
that connects the graph prediction to tumor response.

\section{Measurement and modeling terminology}

Several measurements recur throughout this update. They are all related, but
they answer different questions about the rollout model.

\paragraph{Functional tumor volume.}
Functional tumor volume (FTV) is the estimated volume of active enhancing tumor
tissue on breast MRI. In this work, it is computed from the predicted or
observed supervoxel feature that tracks active tumor volume. FTV is important
because it is the patient-level imaging endpoint that most directly represents
active tumor burden during neoadjuvant treatment. A model can produce a
reasonable-looking tumor cloud while still predicting too little or too much
total active tumor. That is why FTV calibration became the central issue after
the Monte Carlo analysis.

\paragraph{Near-zero T3 MRI FTV diagnostic.}
The rollout model predicts an imaging trajectory through the available MRI time
points. The Monte Carlo script reported
\[
P(\text{FTV}_{T3} < 0.1\text{ mL}),
\]
which answers one imaging question: how often does the simulated T3 MRI tumor
burden become essentially zero under this threshold? I treat this only as a
descriptive MRI FTV diagnostic. It is not used as a training target, loss term,
or primary endpoint in the retraining analysis.

\paragraph{Supervoxels, alive probability, and alive mass.}
The tumor graph is made of supervoxels, which are small 3D regions of tumor
tissue. Each supervoxel has a predicted alive probability at the next visit.
Alive mass is the sum of those probabilities across the graph:
\[
\text{alive mass} = \sum_i P(\text{node } i \text{ alive}).
\]
The observed alive count is the number of supervoxels that are actually present
or active at the observed future visit. Alive mass matters because predicted
FTV is not just the sum of all predicted node volumes; it should be the sum of
active volume from nodes that remain alive.

\paragraph{FTV residuals and log-volume residuals.}
The conditional Monte Carlo sampler is residual based. An aggregate FTV
residual is the difference between observed FTV and deterministic predicted
FTV:
\[
\text{FTV residual} = \text{observed FTV} - \text{predicted FTV}.
\]
This residual adjusts the sampled patient-level tumor burden. Log-volume
residuals are node-level residuals computed in log space:
\[
\log(1 + \text{observed node volume}) -
\log(1 + \text{predicted node volume}).
\]
Log space is used because tumor volumes have a long right tail; it avoids
letting a small number of very large nodes dominate the residual distribution.

\paragraph{Spatial residuals and cloud metrics.}
The sampler also uses spatial residuals. A global displacement residual shifts
the whole predicted tumor cloud coherently, while local displacement residuals
add node-level spatial variation. Spatial quality is evaluated with sliced
Wasserstein distance (SWD), Chamfer distance, and Dice. SWD and Chamfer measure
how far the predicted and observed point clouds are from each other; lower is
better. Dice measures overlap after voxelizing the clouds; higher is better.

\paragraph{Uncertainty metrics.}
Coverage measures how often the observed FTV lands inside the model's predicted
uncertainty interval. Raw 90\% coverage uses the Monte Carlo FTV distribution
directly: the lower endpoint is the 5th percentile and the upper endpoint is the
95th percentile. If the model is well calibrated, about 90\% of held-out
observations should fall inside those intervals. Conformal 90\% coverage starts
from the same raw Monte Carlo interval but then expands it using held-out
nonconformity scores, which are the observed misses between predicted intervals
and true FTV values. This gives an empirical correction when the raw Monte Carlo
intervals are too narrow or systematically off center. Interval width measures
the size of those uncertainty intervals, so it captures sharpness: narrower
intervals are better only when coverage remains adequate. The Continuous Ranked
Probability Score (CRPS) evaluates the entire predicted FTV distribution against
the observed FTV, not just one interval or one point estimate. Lower CRPS means
the distribution puts more probability mass near the observed value while still
remaining appropriately spread out.

\section{Conditional Monte Carlo experiment design}

The Monte Carlo experiment was designed as a stress test of whether a
deterministic graph rollout could be used as the center of a clinically useful
uncertainty distribution. The deterministic model gives one predicted future
tumor graph. The Monte Carlo layer asks a different question: given the errors
seen in held-out patients, what range of future FTV values and tumor clouds are
plausible for a new held-out patient?

\paragraph{Base model and held-out protocol.}
The original experiment used the retained scheduled-sampling consistent rollout
checkpoints:
\[
\texttt{runs/consistent\_forecaster\_v2/v2\_sched\_samp}.
\]
Each patient was evaluated with the held-out fold checkpoint for that patient,
so the checkpoint used for a patient was not trained on that patient. The full
run included 758 patients, 4,548 rollout records, and 1,164,288 Monte Carlo
samples, with 256 samples per rollout record.

The same accounting was used in the retrained follow-up. Each retrained model
produced 4,548 rollout records and 1,164,288 MC samples. Because I ran three
retrained checkpoints, the follow-up generated 3,492,864 retrained MC samples.
Including the original baseline run used for comparison, the analysis compared
4,657,152 MC samples across the baseline and retrained models. The sample count
comes from
\[
758\text{ patients} \times 6\text{ conditioning buckets} \times 256
\text{ MC draws} = 1{,}164{,}288\text{ samples per model}.
\]

\paragraph{Conditioning design.}
The experiment was conditional because it treated each patient's already observed
visits as known information, then sampled uncertainty only over the future
visits. This matches the clinical use case: at different points during
treatment, different amounts of imaging history are available.
The experiment evaluated three conditioning regimes:
\begin{itemize}
    \item T0 observed, then predict T1, T2, and T3.
    \item T0 and T1 observed, then predict T2 and T3.
    \item T0, T1, and T2 observed, then predict T3.
\end{itemize}
This produced six start-visit/predicted-visit buckets:
T0$\rightarrow$T1, T0$\rightarrow$T2, T0$\rightarrow$T3,
T1$\rightarrow$T2, T1$\rightarrow$T3, and T2$\rightarrow$T3. The T3 buckets
were the most important because T3 is the final pre-surgery MRI time point.
Because every patient contributes all six buckets, each model has
$758 \times 6 = 4{,}548$ patient/bucket rollout records before MC sampling.

\paragraph{Calibration residual design.}
The sampler was residual based. First, the deterministic rollout was run for
each patient and each conditioning bucket. Then the difference between the
deterministic prediction and the observed held-out outcome was stored as a
calibration residual. Residuals were bucketed by start visit and predicted visit,
because the error distribution is not the same for T0$\rightarrow$T1 as it is
for T0$\rightarrow$T3. For each target patient, that patient's own residuals
were excluded from the calibration pool. This made the sampling leave-one-patient
out with respect to residual information: the patient being evaluated could not
contribute its own error back into its uncertainty distribution.

For a given target patient and bucket, the patient-level residual candidates are
therefore the other 757 patients in that same bucket. This exclusion is
important because it prevents a patient from borrowing its own observed error
when constructing its MC distribution.

\paragraph{What one Monte Carlo draw did.}
Each Monte Carlo draw started from the deterministic held-out-fold rollout for
one target patient and one conditioning bucket. The model weights, observed
conditioning visits, fold assignment, and deterministic rollout center were not
randomized. What changed across draws was the residual information resampled
from other patients in the same bucket. Each draw randomized:
\begin{enumerate}
    \item one patient-level calibration residual, sampled with replacement from
    the non-target patients in the same start-visit/predicted-visit bucket;
    \item the aggregate FTV residual from that sampled calibration patient,
    defined as observed FTV minus deterministic predicted FTV;
    \item the alive-count residual from that sampled calibration patient,
    defined as observed alive count minus deterministic alive mass;
    \item the global displacement residual from that sampled calibration
    patient, which shifts the predicted tumor cloud coherently;
    \item an alive mask, sampled by Gumbel top-k from the target patient's
    predicted alive probabilities after the sampled alive-count residual sets
    the total number of alive nodes;
    \item local displacement residuals, sampled independently with replacement
    for each node from the pooled node-level residuals in the same bucket; and
    \item log-volume residuals, sampled independently with replacement for each
    node from the pooled node-level volume residuals in the same bucket.
\end{enumerate}

For each draw, the sampler perturbed positions by adding the sampled global
shift plus node-level local displacement noise. It perturbed node volumes in
log space by applying sampled log-volume residuals, then clipped sampled volumes
to be nonnegative. Sampled FTV was computed as the sum of sampled volumes over
the sampled alive nodes, plus the sampled aggregate FTV residual, clipped at
zero. Repeating this 256 times gave an empirical distribution of possible
future FTV values and tumor-cloud summaries for that patient/bucket rather than
a single deterministic rollout.

\paragraph{Why the sampler was held fixed.}
The purpose of the first Monte Carlo experiment was diagnostic, not to tune the
sampler until the final numbers looked better. I therefore kept the sampling
procedure fixed and asked whether the uncertainty distribution exposed a
weakness in the deterministic rollout center. In the retrained follow-up, I
kept that same sampler fixed again and swapped only the base model. This
isolated the question I cared about: if the deterministic FTV center is
corrected, do the Monte Carlo mean, intervals, coverage, CRPS, and spatial
sample summaries become better calibrated without changing the uncertainty
machinery itself?

\paragraph{Outputs and evaluation.}
For every rollout record, the experiment saved the deterministic prediction, the
Monte Carlo sample distribution, per-patient summaries, and calibration
summaries. The main FTV outputs were deterministic bias and MAE, Monte Carlo
mean and median bias/MAE, raw 90\% coverage, conformal 90\% coverage, interval
width, and Continuous Ranked Probability Score (CRPS). The main graph outputs
were alive-count error, sliced Wasserstein distance (SWD), Chamfer distance, and
Dice. This combination was important because a model can have acceptable-looking
spatial rollouts while still being biologically miscalibrated in total tumor
burden, or it can improve FTV while degrading the tumor cloud geometry.

\section{What the baseline Monte Carlo simulation revealed}

This section refers to the original conditional MC run on the baseline
\texttt{v2\_sched\_samp} scheduled-sampling model, before biological
retraining. It is the diagnostic experiment that motivated the retraining, not
the final retrained MC comparison. The first important result was systematic
deterministic FTV underprediction. The T0-to-T3 endpoint was the most important
case:
\[
\text{observed mean T3 FTV} = 24.94\text{ mL}, \qquad
\text{deterministic mean T3 FTV} = 12.43\text{ mL}.
\]
That is a mean signed error of -12.51 mL and an FTV MAE of 12.88 mL. In
biological terms, the model learned a rollout trajectory that shrank active
tumor burden too aggressively.

For the baseline model, the residual Monte Carlo correction did not fix the
problem as a point predictor. It corrected the direction of the bias but
overcorrected the
magnitude. For T0-to-T3, the Monte Carlo mean was 31.30 mL against the observed
mean of 24.94 mL, giving a +6.36 mL bias and a 15.26 mL MC-mean MAE. That told
me the uncertainty layer was partly compensating for a biased base model rather
than expressing uncertainty around a well-centered deterministic prediction.

\begin{figure}[H]
    \centering
    \includegraphics[width=\linewidth]{mc_diagnostic_bias_coverage_width.png}
    \caption{Conditional Monte Carlo diagnostics for the original
    scheduled-sampling rollout model. The deterministic rollout was biased low
    for FTV, while the residual Monte Carlo mean overcorrected high. Raw 90\%
    FTV intervals under-covered, and split-conformal expansion restored nominal
    coverage only with wide intervals.}
    \label{fig:mc-diagnostics}
\end{figure}

\begin{table}[H]
\centering
\small
\caption{Summary of the original conditional Monte Carlo FTV results. Coverage
is for 90\% FTV intervals.}
\label{tab:mc-summary}
\begin{tabular}{lrrrrrr}
\toprule
Bucket & Obs mean & Det mean & MC mean & Det bias & MC bias & Raw / Conf cov. \\
\midrule
T0$\rightarrow$T1 & 24.31 & 12.61 & 30.26 & -11.70 & 5.95 & 0.752 / 0.901 \\
T0$\rightarrow$T2 & 24.75 & 12.16 & 31.30 & -12.58 & 6.56 & 0.793 / 0.901 \\
T0$\rightarrow$T3 & 24.94 & 12.43 & 31.30 & -12.51 & 6.36 & 0.766 / 0.901 \\
T1$\rightarrow$T2 & 24.75 & 12.18 & 31.61 & -12.56 & 6.86 & 0.793 / 0.901 \\
T1$\rightarrow$T3 & 24.94 & 12.20 & 31.85 & -12.74 & 6.90 & 0.764 / 0.901 \\
T2$\rightarrow$T3 & 24.94 & 12.19 & 31.97 & -12.75 & 7.02 & 0.776 / 0.901 \\
\bottomrule
\end{tabular}
\end{table}

The second baseline MC result was that later conditioning helped volume scoring,
but did not fully solve spatial uncertainty. T3 FTV CRPS improved as more
observed visits were included, from 8.35 for T0-to-T3 to 7.49 for
T0,T1,T2-to-T3. However, the corresponding SWD worsened from 1.32 mm to
2.97 mm in the Monte Carlo summaries. This suggested that the sampler was
capturing some volume behavior but was still weak as a spatial uncertainty
model.

\begin{figure}[H]
    \centering
    \includegraphics[width=0.78\linewidth]{mc_t3_conditioning_crps_swd.png}
    \caption{T3 conditional Monte Carlo behavior. Later conditioning improved
    FTV CRPS, but SWD worsened, showing that volume uncertainty and spatial
    uncertainty were not improving together.}
    \label{fig:mc-conditioning}
\end{figure}

The third baseline MC result was about the near-zero T3 MRI FTV diagnostic. The
field was computed as
\[
P(\text{FTV}_{T3} < 0.1\text{ mL}),
\]
which only asks whether the simulated T3 MRI FTV is essentially zero under a
fixed threshold. In this run the value was almost always near zero, around
0.007 to 0.009 at T3. This means the threshold is very strict and should be
treated as a descriptive imaging diagnostic, not as a central response endpoint.
The threshold was not part of the retraining objective; the retrained models
learned FTV and alive-burden targets.

The fourth baseline MC result was alive-count miscalibration. The original
conditional MC experiment had T3 alive-count absolute errors around 39 to 42
supervoxels. Since alive mass controls which nodes contribute to active tumor
volume, this was not a cosmetic issue; it directly affected FTV calibration.

\section{Why further exploration was justified}

The Monte Carlo experiment changed the interpretation of the earlier success.
The model was good enough to demonstrate coherent graph rollout behavior, and
the conformal wrapper could recover nominal interval coverage. However, the
center of the uncertainty distribution was biologically miscalibrated.

The main lesson was:
\[
\text{good spatial rollout metrics do not guarantee calibrated tumor burden.}
\]
The model had learned local graph dynamics, but the clinically important
endpoint was not directly controlled. If the deterministic rollout is wrong by
about 12 to 13 mL at T3, downstream Monte Carlo response summaries inherit that
bias. This motivated a targeted follow-up question:

\begin{quote}
Can I keep the same consistent graph rollout structure, but add training
pressure so the rollout is also calibrated for FTV and alive tumor burden?
\end{quote}

\section{Biological retraining response}

The retraining grid kept the same consistent graph rollout model family, but
added biological endpoint supervision while retaining the scheduled-sampling
rollout setup:
\begin{itemize}
    \item aggregate FTV loss, penalizing patient-level error in total active
    tumor burden;
    \item alive supervision, encouraging correct alive/dead supervoxel
    behavior and alive mass; and
    \item the same held-out fold protocol used for the original
    scheduled-sampling model.
\end{itemize}

This distinction matters. The MC results reported below evaluate the bio
retraining checkpoints. Those checkpoints added FTV and alive-burden objectives,
but they were not the later random-start-plus-horizon rerun. Therefore, this
analysis isolates whether biological endpoint calibration alone fixes the MC
calibration problem before changing the rollout-start protocol.

The evaluated variants were:
\[
\begin{array}{ll}
\texttt{bio\_ftv000\_alive002} & \text{alive-only biological retraining},\\
\texttt{bio\_ftv005\_alive001} & \text{low FTV + low alive},\\
\texttt{bio\_ftv010\_alive002} & \text{medium FTV + medium alive},\\
\texttt{bio\_ftv020\_alive005} & \text{high FTV + high alive},\\
\texttt{bio\_ftv010\_alive000} & \text{FTV-only ablation}.\\
\end{array}
\]

The retraining results were much better aligned with the Monte Carlo diagnosis.
The baseline T0-to-T3 FTV MAE was 12.88 mL. The best retrained candidate,
\texttt{bio\_ftv020\_alive005}, reduced this to 5.88 mL, a 54.3\% reduction.
Its mean T3 FTV prediction was 23.21 mL against an observed mean of 24.94 mL,
leaving a mean signed error of -1.73 mL rather than -12.51 mL.

\begin{figure}[H]
    \centering
    \includegraphics[width=\linewidth]{bio_retraining_ftv_mae_bias.png}
    \caption{FTV endpoint calibration after biological retraining. The
    retrained candidates reduced T0-to-T3 FTV MAE by roughly half and removed
    most of the systematic underprediction seen in the baseline rollout model.}
    \label{fig:bio-ftv}
\end{figure}

\begin{table}[H]
\centering
\small
\caption{T0-to-T3 evaluation of the baseline and biological retraining
candidates. Lower FTV MAE is better.}
\label{tab:bio-summary}
\begin{tabular}{lrrrrr}
\toprule
Model & Pred FTV & Obs FTV & Bias & FTV MAE & MAE reduction \\
\midrule
\texttt{v2\_sched\_samp} & 12.43 & 24.94 & -12.51 & 12.88 & 0.0\% \\
\texttt{bio\_ftv000\_alive002} & 24.40 & 24.94 & -0.54 & 6.39 & 50.4\% \\
\texttt{bio\_ftv005\_alive001} & 23.85 & 24.94 & -1.10 & 6.15 & 52.2\% \\
\texttt{bio\_ftv010\_alive002} & 23.44 & 24.94 & -1.50 & 5.99 & 53.5\% \\
\texttt{bio\_ftv020\_alive005} & 23.21 & 24.94 & -1.73 & 5.88 & 54.3\% \\
\texttt{bio\_ftv010\_alive000} & 23.44 & 24.94 & -1.51 & 5.99 & 53.5\% \\
\bottomrule
\end{tabular}
\end{table}

The paired patient-level comparison gave the same conclusion. For
\texttt{bio\_ftv020\_alive005}, the mean paired FTV MAE improvement was
7.00 mL, with an approximate 95\% interval of 6.06 to 7.93 mL, and 79.6\% of
patients improved.

\begin{figure}[H]
    \centering
    \includegraphics[width=0.82\linewidth]{bio_retraining_paired_ftv_improvement.png}
    \caption{Patient-level paired FTV MAE improvement relative to the baseline
    scheduled-sampling rollout model. Positive values mean the retrained model
    reduced FTV absolute error for the same patient.}
    \label{fig:paired-improvement}
\end{figure}

\section{Geometry was preserved}

The important caveat is that the retrained models did not simply become crude
volume predictors. The spatial metrics stayed essentially flat. For the best
FTV candidate, mean SWD changed by about 0.003 mm and Chamfer changed by about
0.006 mm relative to the baseline. Dice changed by less than 0.001.

\begin{figure}[H]
    \centering
    \includegraphics[width=\linewidth]{bio_retraining_geometry_deltas.png}
    \caption{Changes in spatial metrics relative to the original
    scheduled-sampling baseline. The FTV improvement was achieved without a
    meaningful collapse in spatial rollout quality.}
    \label{fig:geometry-deltas}
\end{figure}

Alive-count error also improved dramatically in the retraining evaluations,
from 58.00 in the baseline simulation evaluation to about 1.70 in the
biological retraining variants.

\begin{figure}[H]
    \centering
    \includegraphics[width=0.82\linewidth]{bio_retraining_alive_error.png}
    \caption{Alive-count absolute error in the baseline and retraining
    evaluations. The improvement is large, but the current ablation grid does
    not isolate the explicit alive-loss term as the sole driver because the
    FTV-only ablation also improved.}
    \label{fig:alive-error}
\end{figure}

The alive term should be interpreted cautiously. The alive-count
improvement was identical across the biological retraining variants, including
the FTV-only ablation. The current evidence strongly supports endpoint-aware
retraining, especially the aggregate FTV objective, but it does not cleanly
prove that explicit alive supervision is the mechanism behind the alive-count
improvement.

\section{Updated conditional Monte Carlo results on retrained models}

After the deterministic retraining analysis, I reran the conditional Monte
Carlo sampler on three FTV-calibrated checkpoints:
\texttt{bio\_ftv020\_alive005}, \texttt{bio\_ftv010\_alive000}, and
\texttt{bio\_ftv010\_alive002}. I kept the sampler unchanged: same residual
buckets by start and predicted visit, same target-patient exclusion from the
calibration pool, same Gumbel top-k alive-mask sampling, same split-conformal
interval calculation, same 758-patient cohort, and 256 MC samples per rollout
record. This was important because the question was not whether a new sampler
could be tuned to look better. The question was whether fixing the deterministic
FTV center would make the existing uncertainty distribution better calibrated.

The answer was yes. For the primary model, \texttt{bio\_ftv020\_alive005},
T0-to-T3 deterministic bias stayed close to zero at -1.73 mL instead of
returning to the old -12.51 mL bias. The MC mean no longer had to compensate for
a badly undercentered base model: T0-to-T3 MC-mean FTV MAE dropped from
15.26 mL to 7.61 mL. Raw 90\% FTV coverage improved from 0.766 to 0.876 before
conformal correction, and the conformal 90\% interval width shrank from
59.51 mL to 27.95 mL. CRPS improved from 8.35 to 5.16, and MC alive-count error
fell from 40.78 to 2.69 supervoxels.

\begin{figure}[H]
    \centering
    \includegraphics[width=\linewidth]{mc_bio_retrained_t0t3_summary.png}
    \caption{T0-to-T3 conditional MC comparison between the original
    scheduled-sampling baseline and the three FTV-calibrated retrained models.
    The retrained models reduce MC-mean FTV error, improve raw 90\% coverage,
    shrink conformal intervals, and greatly reduce alive-count error.}
    \label{fig:mc-bio-t0t3}
\end{figure}

\begin{table}[H]
\centering
\small
\caption{T0-to-T3 conditional MC comparison. Coverage is for raw 90\% FTV
intervals before conformal expansion. Width is the conformal 90\% FTV interval
width.}
\label{tab:mc-bio-t0t3}
\begin{tabular}{lrrrrrrr}
\toprule
Model & Det bias & Det MAE & MC bias & MC MAE & Raw cov. & Width & Alive err. \\
\midrule
\texttt{v2\_sched\_samp} & -12.51 & 12.88 & 6.36 & 15.26 & 0.766 & 59.51 & 40.78 \\
\texttt{bio\_ftv020\_alive005} & -1.73 & 5.88 & 1.93 & 7.61 & 0.876 & 27.95 & 2.69 \\
\texttt{bio\_ftv010\_alive000} & -1.51 & 5.99 & 1.14 & 7.31 & 0.875 & 26.92 & 7.00 \\
\texttt{bio\_ftv010\_alive002} & -1.50 & 5.99 & 1.76 & 7.53 & 0.883 & 29.14 & 2.70 \\
\bottomrule
\end{tabular}
\end{table}

The improvement also held across the three T3 conditioning buckets. For
\texttt{bio\_ftv020\_alive005}, MC-mean FTV MAE was 7.61 mL for T0-to-T3,
7.18 mL for T1-to-T3, and 5.80 mL for T2-to-T3. The corresponding conformal
widths were 27.95 mL, 29.61 mL, and 24.19 mL, all far below the old baseline
range of about 59 to 60 mL. This matters because the MC improvement was not
only a single T0-to-T3 artifact; it persisted when later observed imaging visits
were used as conditioning information.

\begin{figure}[H]
    \centering
    \includegraphics[width=\linewidth]{mc_bio_retrained_t3_buckets.png}
    \caption{T3 conditional MC metrics across available imaging history. The
    FTV-calibrated retrained models reduce MC-mean FTV MAE, conformal interval
    width, and CRPS across T0-to-T3, T1-to-T3, and T2-to-T3 conditioning.}
    \label{fig:mc-bio-t3-buckets}
\end{figure}

\subsection{Patient-level uncertainty visualization}

The aggregate metrics are necessary, but they are not enough by themselves. The
new patient-level plots show what the MC distribution is doing for individual
rollouts. In these plots, the black curve is the observed FTV trajectory, the
gray curve and band show the original baseline MC mean plus or minus one
standard deviation, and the green curve and band show the retrained primary
model. These are MRI FTV uncertainty trajectories from the T0-conditioned
rollout.

\begin{figure}[H]
    \centering
    \includegraphics[width=\linewidth]{mc_patient_trajectory_mean_std_examples.png}
    \caption{Example patient-level T0-to-T3 MC trajectories. The retrained
    primary model generally shifts the MC mean closer to the observed FTV path
    and reduces the width of the one-standard-deviation band. The examples were
    selected to show low, median, high, and strongly improved T3 cases, so the
    figure is illustrative rather than a replacement for cohort-level
    calibration metrics.}
    \label{fig:mc-patient-trajectories}
\end{figure}

The cohort-wide sorted plot gives the same comparison for all 758 patients in
the T0-to-T3 setting. Patients are ordered by observed T3 FTV, so this figure
shows whether the MC distribution follows the empirical endpoint distribution
across low-, moderate-, and high-volume cases. The baseline MC distribution has
a visibly wide band and often sits high after residual correction. The retrained
model is more centered and sharper. That interpretation is supported by the
quantitative results: raw coverage improved from 0.766 to 0.876 while the
conformal width shrank from 59.51 mL to 27.95 mL. In other words, the narrower
band is not just overconfidence; it is paired with better empirical coverage
and lower CRPS.

\begin{figure}[H]
    \centering
    \includegraphics[width=\linewidth]{mc_t0t3_patient_uncertainty_sorted.png}
    \caption{Cohort-wide T0-to-T3 MC uncertainty bands, with patients sorted by
    observed T3 FTV. The retrained primary model has a narrower and better
    centered MC distribution across the cohort, while the baseline distribution
    remains broad and shifted upward after residual correction.}
    \label{fig:mc-sorted-patient-uncertainty}
\end{figure}

The standard-deviation summary isolates uncertainty width directly. From T0
conditioning, the baseline mean MC FTV standard deviation increased from
24.71 mL at T1 to 27.56 mL at T3. The retrained primary model had substantially
smaller standard deviations: 12.60 mL at T1, 15.33 mL at T2, and 16.10 mL at
T3. This is exactly the expected pattern if the old MC layer was compensating
for a biased deterministic center. Once the deterministic center was corrected,
the residual distribution no longer had to spread as widely to cover the
observed endpoint.

\begin{figure}[H]
    \centering
    \includegraphics[width=0.88\linewidth]{mc_t0_to_t3_std_distribution.png}
    \caption{Patient-level MC FTV standard deviation by predicted visit from
    T0 conditioning. The retrained primary model has consistently lower
    patient-level uncertainty width than the original baseline, especially at
    the final T3 endpoint.}
    \label{fig:mc-std-distribution}
\end{figure}

All three retrained candidates passed the prespecified T0-to-T3 decision gates.
The FTV-only ablation had the lowest T0-to-T3 MC-mean FTV MAE, 7.31 mL, and the
narrowest conformal interval width, 26.92 mL. However, its MC alive-count error
was 7.00 supervoxels, compared with 2.69 for
\texttt{bio\_ftv020\_alive005} and 2.70 for
\texttt{bio\_ftv010\_alive002}. I therefore interpret the FTV-only ablation as
evidence that aggregate FTV supervision explains much of the volume calibration
gain, but I would still keep \texttt{bio\_ftv020\_alive005} as the primary
balanced candidate because it combines strong FTV calibration with much better
alive-count calibration.

The main residual caution is subgroup drift. For the primary model at
T0-to-T3, the HR-/HER2+ subgroup remained the hardest subgroup: deterministic
FTV MAE was 8.65 mL, MC-mean FTV MAE was 10.78 mL, raw coverage was 0.842, and
CRPS was 8.07. The near-zero T3 MRI FTV diagnostic remains an
imaging-threshold summary, not a primary endpoint in this analysis.

\section{Progress interpretation}

The concise interpretation is:

\begin{quote}
I had a model that worked well as a rollout model, but when I used it as the
center of a conditional Monte Carlo simulator, I found a systematic biological
calibration problem. It underpredicted T3 FTV by about 12.5 mL. The residual MC
layer could recover nominal coverage after conformal expansion, but the
intervals were wide and the MC mean overcorrected. That told me the uncertainty
method was compensating for a biased base model.

Because FTV is the key imaging endpoint for tumor burden, I retrained the same
consistent graph model with aggregate FTV and alive-burden losses. The best
retrained model cut T0-to-T3 FTV MAE from 12.88 mL to 5.88 mL while preserving
the spatial metrics. When I reran the same conditional MC sampler on the
retrained models, the uncertainty distribution also improved: the primary model
cut T0-to-T3 MC-mean FTV MAE from 15.26 mL to 7.61 mL, increased raw 90\%
coverage from 0.766 to 0.876, and reduced conformal interval width from
59.51 mL to 27.95 mL. The follow-up was therefore not a random model tweak; it
was a targeted response to a failure mode revealed by the Monte Carlo
simulation, and the completed MC rerun shows that correcting the deterministic
FTV center materially improved uncertainty calibration.
\end{quote}

\section{Current conclusion and next step}

The current conclusion is that \texttt{bio\_ftv020\_alive005} is the best
balanced candidate from this phase. The FTV-only ablation slightly wins on
T0-to-T3 MC-mean FTV MAE and conformal width, but the high-FTV/high-alive model
has much better alive-count calibration while still passing every FTV and
uncertainty gate. That balance matters because the model is supposed to remain a
graph rollout model, not only a scalar FTV predictor. The patient-level
trajectory plots make the same point visually: the retrained model produces a
more credible FTV uncertainty distribution because it is simultaneously closer
to the observed trajectory, narrower, better covered, and lower by CRPS.

The next methodological decision is whether the corrected random-start-plus-
horizon model family should replace this scheduled-sampling bio grid as the
final architecture. If so, the same MC experiment should be repeated on that
checkpoint family, using the same cohort and sampler, so that the comparison is
not confounded by changes in the uncertainty machinery. Separately, the
subgroup results should be kept visible: HR-/HER2+ patients remain the hardest
group in the primary T0-to-T3 MC run, so that subgroup needs targeted error
review before making strong claims about uniform calibration.

The current MC output supports calibrated T3 MRI FTV uncertainty. Any separate
clinical outcome model should be treated as a distinct prediction problem built
from appropriate outcome labels and covariates, rather than folded into this
rollout uncertainty analysis.

\end{document}
